\documentclass[../main.tex]{subfiles}
\begin{document}
%%%%%%%%%%%%%%%%%%%%%%%%%%%%%%%%%%%%%%%%%%%%%%%%%%%%%%%%%%%%%%%%%%%%%%%%%%%%%%%%%%%%%%%%%%%%%%%%%%%
\chapter{Using and Creating Custom Workspaces}

Workspaces enable you to define how you work in OmniGraffle. With the latest update to
OmniGraffle, we have taken Workspaces a bit further to help you not just define where and how
the inspectors appear, but to customize them for how you work.

\section{Using Workspaces}

As you learned in the previous chapter, OmniGraffle's inspectors are located to the right of the
canvas. This is known as the \textbf{Sidebar Workspace}.

\ogfigure{0.800}{res/og7_workspaces_hero.png}

In addition to the Sidebar workspace, OmniGraffle has two other Workspace states for the
inspectors: \textbf{Floating} and \textbf{Palettes} Workspaces. Each Workspace is assigned a
\textit{function} key as their keyboard shortcut:

\begin{itemize}
  \item Sidebar: \textbf{F2}
  \item Floating: \textbf{F3}
  \item Palettes: \textbf{F4}
\end{itemize}

If you are using a MacBook Pro, MacBook Air, or an Apple
\href{http://www.apple.com/shop/product/MLA22LL/A/magic-keyboard-us-english?fnode=56}{Magic
Keyboard} with your Mac, you also need to press the \textbf{fn} key---located at the
lower-left corner of your keyboard---along with the applicable function key. For example,
\textbf{fn-F4} to open the Palettes workspace. If you have an
\href{http://www.apple.com/shop/product/MB110LL/B/apple-keyboard-with-numeric-keypad-us-english?fnode=56}{Apple
Keyboard with Numeric Keypad}, the \textbf{fn} key is located below the \textbf{F13} key.

Subsequent Workspaces that you add to OmniGraffle are assigned their own function key shortcut.

\begin{ogtip}
  Having the Inspectors available as a floating window or as individual floating palettes allows
  you to take advantage of big external displays for OmniGraffle's main window while placing the
  inspectors off to the side or on a secondary display (such as a laptop screen or another large
  display).
\end{ogtip}

Use \textbf{Command--1} through \textbf{Command--4} to switch between the Object, Properties,
Canvas, and Document inspector tabs (or windows, if you're using the Palettes Workspace),
respectively.

\subsection{Using the Floating Workspace}

To open the Floating Workspace, choose \textbf{Inspectors} $\blacktriangleright$
\textbf{Floating}, or press \textbf{F3}. The Floating Workspace contains all of the tabs you
normally find in Inspector Sidebar, and you can switch between inspector sets using the
default key commands, noted earlier.

\ogfigure{0.800}{res/og7_workspaces01.png}

\subsection{Using the Inspector Palettes Workspace}

To open the inspector palettes, choose \textbf{Inspectors} $\blacktriangleright$
\textbf{Palettes}, or press \textbf{F4}. You'll see a stack of windows appear over
OmniGraffle's main window. You'll need to move them aside individually, but once you've got
them set up---especially if you have them all on a second display---you'll be amazed at how
addicted you can get to having the inspectors in palette form.

\ogfigure{0.800}{res/og7_workspaces02.png}

\section{Creating Custom Workspaces}

In addition to the default Workspaces available from the \textbf{Inspectors} menu, you can also
create custom Workspaces. You can also use Workspaces to save default \textit{states} of
certain inspectors, or a totally custom Workspace window that you tailor to suit your
workstyle.

\subsection{Saving Inspector States as a Workspace}

If you frequently need access to only the Fill and Stroke panes (for example), follow these
steps to create a custom Workspace:

\begin{enumerate}
  \item Choose \textbf{Inspectors} $\blacktriangleright$ \textbf{Object} (\textbf{Command-1}) to
    view the Object inspectors.
  \item \textbf{Option-click} the \textbf{Fill} inspector; this opens the Fill inspector and
    collapses the others in the sidebar.
  \item Click the disclosure triangle next to the \textbf{Stroke} inspector to open that
    inspector as well.
  \item Choose \textbf{Inspectors} $\blacktriangleright$ \textbf{Edit Workspaces}; the following
    window appears:

    \ogfigure{0.800}{res/og7_config_workspaces01.png}
  \item At the bottom of the Workspaces window, click \textbf{+} (for Add). You will hear the
    shutter sound associated with taking a picture; this is OmniGraffle's way of letting you
    know that it has captured the state of the inspector window.
  \item Give this new workspace a name, such as \textbf{Fill-Stroke}.

    \ogfigure{0.800}{res/og7_config_workspaces02.png}
  \item Close the Workspaces window.
\end{enumerate}

Now, whenever you need to just tweak an object's Fill or Stroke, you can press \textbf{F5} to
gain quick access to these two inspector panes.

\subsection{Creating a Custom Workspace from Various Inspectors}

One little-known feature of OmniGraffle is that you can \textit{tear away} inspectors from the
Inspector sidebar. For example, press \textbf{Command-1} to view the Object inspector tab,
click on the handle along the right side of the Shape inspector, and then drag the inspector
to the canvas.

\ogfigure{0.800}{res/og7_config_workspaces03.png}

As you drag the Shape inspector away from the sidebar, you'll notice that it disappears from the
sidebar. When you let go of the mouse, though, the Shape inspector appears in its own floating
window.

\ogfigure{0.769}{res/og7_config_workspaces04.png}

You can switch to the other inspector tabs and drag additional inspectors into the window; for
example, do the following:

\begin{itemize}
  \item From the Object inspectors, drag in the \textbf{Shadow} and \textbf{Geometry} inspectors
  \item From the Properties inspectors (\textbf{Command-2}), drag in the \textbf{Connections}
    inspector
\end{itemize}

As you drag the various inspectors from the Inspector sidebar into the window, you'll notice
that they also disappear entirely from the sidebar. That's okay, and we'll show you how to get
those inspectors back into their respective places in the Inspector sidebar later.

The floating inspector window should appear as follows:

\ogfigure{0.769}{res/og7_config_workspaces05.png}

As you can see, the window's titlebar shows its name as \textbf{Custom 1}; however, you can save
this configuration of inspectors as a Workspace by doing the following:

\begin{enumerate}
  \item Choose \textbf{Inspectors} $\blacktriangleright$ \textbf{Edit Workspaces} to open the
    Workspaces window:

    \ogfigure{0.800}{res/og7_config_workspaces06.png}
  \item Click the \textbf{+} (Add) button, and replace the \textbf{Untitled} label\ldots{}

    \ogfigure{0.800}{res/og7_config_workspaces07.png}

    \ldots{}with something more specific to the inspectors you've chosen; for example,
    \textbf{Diagramming}:

    \ogfigure{0.800}{res/og7_config_workspaces08.png}
  \item Close the Workspaces window.
\end{enumerate}

As you may have noticed, whenever you create a new Workspace, OmniGraffle assigns one of the
Function keys as a shortcut for opening that particular Workspace. However, the keyboard
shortcut doesn't close the Workspace window. Instead, you will need to click the red close
button in the Workspace window's titlebar.

When you close the Workspace window, you'll notice that its individual inspectors do not return
to the Inspector sidebar. To get the inspectors back in the sidebar, you can either choose
\textbf{Window} $\blacktriangleright$ \textbf{Sidebar} or press \textbf{F2}. If you were to
press \textbf{F5} to open the Diagramming Workspace, the inspectors pop out of the sidebar
again and take up residence in the \textbf{Custom 1} window until you return them to the
sidebar.
%%%%%%%%%%%%%%%%%%%%%%%%%%%%%%%%%%%%%%%%%%%%%%%%%%%%%%%%%%%%%%%%%%%%%%%%%%%%%%%%%%%%%%%%%%%%%%%%%%%
\end{document}
